\chapter{Аналитический обзор и теоретическая база}
\label{ch:overview}

\section{Постановка задачи системного анализа вычислительной платформы}

Под системным анализом вычислительной платформы в рамках данной работы понимается совокупность методов, моделей и практических процедур, направленных на получение воспроизводимого, интерпретируемого и прикладно значимого представления о состоянии аппаратно-программного комплекса и о его предельных возможностях при различных профилях нагрузки. В отличие от локального профилирования отдельного приложения, где исследователь сосредоточен на поведении конкретного алгоритма или программного модуля, системный анализ охватывает платформу в целом. Это означает необходимость рассматривать не только центральный процессор, но и память, дисковую и сетевую подсистемы, а также влияние операционной системы, фоновых процессов, механизмов энергосбережения и особенностей используемой программной среды.

Необходимость такого подхода обусловлена тем, что современная вычислительная система представляет собой сложную многоуровневую структуру, в которой итоговая производительность определяется не отдельным компонентом, а согласованностью взаимодействия всех подсистем. Даже высокопроизводительный процессор не гарантирует высокой прикладной эффективности, если приложение ограничено пропускной способностью памяти, скоростью накопителя, сетевой задержкой или особенностями планировщика. Аналогично, высокие паспортные характеристики оборудования не всегда отражают поведение системы в реальных условиях эксплуатации, где присутствуют фоновые службы, антивирусная активность, индексация файлов, обновления, виртуализация и иные процессы, способные существенно изменить измеряемый результат.

В теоретическом плане задача системного анализа тесно связана с понятием производительности как многомерной характеристики. Под производительностью нельзя понимать исключительно скорость выполнения одной инструкции, одного потока или одной операции. Для прикладного анализа существенны по меньшей мере следующие аспекты: вычислительная мощность процессора, способность платформы эффективно использовать параллелизм, скорость и латентность обращения к памяти, интенсивность и устойчивость ввода-вывода, а также производительность криптографических примитивов, если система применяется в задачах защиты информации. Следовательно, корректный анализ должен учитывать разнообразие классов нагрузок и невозможность свести состояние платформы к одному универсальному числу.

С практической точки зрения системный анализ востребован в ряде типовых сценариев. Во-первых, он необходим для первичной диагностики рабочих станций и серверов, когда требуется установить, какие именно ресурсы ограничивают производительность прикладных программ. Во-вторых, он важен для сравнения аппаратных платформ при выборе оборудования под определённый класс задач, например под компиляцию, обработку больших массивов данных, криптографические сервисы или нагрузку, чувствительную к памяти. В-третьих, системный анализ имеет значимость в образовательном и исследовательском процессе, поскольку позволяет наглядно показать, каким образом архитектурные особенности вычислительной системы проявляются в измеряемых показателях.

Отдельно следует отметить, что в условиях современных гетерогенных и многоядерных архитектур интерпретация результатов становится значительно сложнее, чем в эпоху относительно однородных процессоров. Наличие производительных и энергоэффективных ядер, многоуровневых кэшей, аппаратных блоков ускорения и развитых механизмов динамического изменения частоты приводит к тому, что поведение системы начинает существенно зависеть от характера нагрузки и политики управления ресурсами. По этой причине результаты, полученные без учёта контекста исполнения, могут быть формально корректными, но недостаточно информативными для практических выводов.

Ещё одной принципиальной сложностью является стохастическая природа наблюдаемых измерений. В многозадачной операционной системе любое тестирование выполняется на фоне иных процессов и событий, часть которых не контролируется исследователем напрямую. На результат способны влиять миграция потоков между ядрами, прерывания от устройств, обращение к файловой системе, состояние кэшей, колебания частоты CPU, прогрев компонентов и сетевые события. Это означает, что системный анализ должен включать не только сбор значений, но и механизмы повышения устойчивости выводов: повторные измерения, предварительный разогрев, фильтрацию выбросов и оценку разброса.

В рамках настоящей работы применяется комбинированный подход к решению данной задачи. С одной стороны, используется мониторинг состояния системы в реальном времени, позволяющий зафиксировать контекст выполнения измерений: загрузку процессора, распределение нагрузки по ядрам, использование памяти, дисковую и сетевую активность, а также присутствие конкурирующих процессов. С другой стороны, используются синтетические тесты, создающие контролируемую нагрузку и позволяющие оценить возможности платформы в воспроизводимой форме. Именно сочетание этих двух компонентов делает анализ более содержательным: мониторинг отвечает на вопрос о текущем состоянии системы, а бенчмарки --- на вопрос о её поведении при специально организованных вычислительных сценариях.

Таким образом, постановка задачи системного анализа вычислительной платформы в данной работе сводится к созданию такого подхода и такого программного инструмента, которые позволяют одновременно наблюдать состояние системы, формировать репрезентативную нагрузку, сопоставлять измерения с фоновыми факторами и получать результаты, пригодные для практической интерпретации. Это и определяет логику дальнейшего исследования, в котором рассматриваются классы метрик, существующие подходы и программные средства, а также формулируются требования к разрабатываемому комплексу.

\section{Классы метрик и источники вариативности измерений}

Для практического анализа вычислительной платформы целесообразно выделять несколько групп метрик, каждая из которых характеризует отдельный аспект функционирования системы. Подобная классификация важна не только для упорядочивания наблюдаемых данных, но и для корректной интерпретации результатов, поскольку разные классы нагрузок опираются на разные аппаратные и программные ресурсы. Следовательно, методика анализа должна учитывать различную природу вычислительных, память-зависимых, дисковых, сетевых и криптографических процессов.

\textbf{Процессорные метрики} отражают способность платформы выполнять вычисления и эффективно использовать имеющийся параллелизм. К данной группе относятся общая загрузка CPU, загрузка по отдельным логическим ядрам, доступные частотные режимы, средняя нагрузка системы, число физических и логических ядер, а также косвенные признаки конкуренции за процессорное время. Однако интерпретация этих данных требует осторожности. Например, высокая суммарная загрузка CPU не всегда означает, что приложение эффективно использует ресурсы: она может быть вызвана конкуренцией нескольких процессов, неудачным распараллеливанием или повышенной активностью фоновых служб. Напротив, умеренная загрузка может сопровождаться реальным дефицитом вычислительных ресурсов, если потоки часто блокируются или ожидают данных из памяти и ввода-вывода. Поэтому при анализе процессорных показателей важно учитывать не только абсолютный уровень загрузки, но и форму распределения нагрузки по ядрам, динамику частоты, устойчивость режима исполнения и характер выполняемой работы.

С инженерной точки зрения процессорные метрики особенно значимы для задач, чувствительных к интенсивности вычислений, ветвлениям и масштабированию по потокам. К ним относятся компиляция, численные методы, обработка мультимедийных данных, моделирование и часть криптографических операций. Для таких сценариев важны не только средние значения загрузки, но и способность системы длительно поддерживать вычислительную активность без деградации из-за троттлинга, ограничений энергопотребления или неэффективного планирования. Именно поэтому в прикладных средствах анализа полезно сопоставлять данные о загрузке ядер с результатами синтетических CPU-бенчмарков.

\textbf{Метрики памяти} образуют вторую принципиально важную группу, поскольку для многих современных приложений именно подсистема памяти становится фактическим ограничителем производительности. К этим метрикам относятся общий объём оперативной памяти, объём доступной и занятой памяти, использование подкачки, косвенные признаки memory pressure, а также показатели, характеризующие качество доступа к данным. В прикладном отношении особенно значимы два параметра: пропускная способность и латентность. Пропускная способность показывает, насколько быстро система способна обрабатывать последовательные потоки данных, тогда как латентность определяет время реакции на нерегулярные обращения к памяти. Для задач потоковой обработки, мультимедиа, больших матричных вычислений и копирования массивов критична именно пропускная способность. Для компиляторов, графовых алгоритмов, баз данных и многих аналитических приложений большое значение имеет латентность случайного доступа.

Особенность памяти как объекта анализа заключается в том, что её влияние на производительность часто носит скрытый характер. Приложение может демонстрировать внешне умеренную загрузку CPU, но фактически быть ограниченным не вычислениями как таковыми, а задержками получения данных из кэша более низкого уровня или из основной памяти. По этой причине измерение только объёма занятой памяти не даёт достаточной информации для системного анализа. Требуются дополнительные тесты, позволяющие оценить скорость линейных операций и чувствительность платформы к нерегулярному шаблону обращений.

\textbf{Метрики подсистемы ввода-вывода} включают показатели работы накопителей и сетевых интерфейсов. Для дисковой подсистемы важны скорости чтения и записи, интенсивность операций, характер доступа к данным, влияние буферизации и наличие конкурирующей фоновой активности. Для сетевой подсистемы существенны скорости передачи и приёма данных, наличие активных соединений и общая интенсивность сетевого обмена. Следует учитывать, что для многих приложений ввод-вывод выступает не просто вспомогательным компонентом, а фактором, определяющим предельную производительность. Это особенно характерно для резервного копирования, работы с большими файлами, систем сборки, обработки журналов, потоковой передачи данных и распределённых сервисов.

Системный анализ ввода-вывода осложняется тем, что наблюдаемые скорости сильно зависят от текущего состояния очередей операций, кэширования со стороны ОС, параллельной активности других программ и особенностей файловой системы. Например, высокие показатели чтения могут объясняться не только быстрым накопителем, но и попаданием данных в системный кэш. Аналогично, всплеск записи может быть связан с временной буферизацией, а не с устойчивой физической скоростью устройства. Следовательно, для интерпретации I/O-метрик необходим контекст в виде дополнительной телеметрии и повторных измерений.

\textbf{Криптографические метрики} имеют особое значение в рассматриваемой работе, поскольку значительная часть современных приложений функционирует в среде, где защита данных является обязательным требованием. К этой группе относятся пропускная способность симметричного шифрования, скорость вычисления хеш-функций, производительность асимметрических операций и эффективность алгоритмов аутентифицированного шифрования. На практике именно эти показатели определяют, насколько быстро система может обслуживать защищённые соединения, выполнять цифровую подпись, проверять целостность файлов и обрабатывать конфиденциальные данные.

Криптографические метрики сложны для анализа по нескольким причинам. Во-первых, разные алгоритмы по своей природе нагружают различные элементы вычислительной платформы: одни зависят преимущественно от арифметической мощности и ширины векторных инструкций, другие --- от производительности операций с большими целыми числами, третьи --- от особенностей доступа к памяти. Во-вторых, результат существенно зависит от того, поддерживает ли платформа специализированные аппаратные инструкции ускорения. В-третьих, сами криптографические библиотеки могут быть реализованы по-разному, что делает интерпретацию абсолютных чисел без знания программного контекста неполной. Тем не менее именно включение криптографических тестов в состав общего системного анализа делает инструмент особенно полезным для задач компьютерной безопасности.

Наряду с перечисленными группами метрик следует учитывать и \textbf{контекстные системные показатели}, которые не являются целевым объектом бенчмаркинга, но оказывают прямое влияние на достоверность выводов. К ним можно отнести время работы системы, наличие активных фоновых процессов, интенсивность обслуживания диска и сети, а также общий уровень загруженности ОС. Эти данные помогают объяснить причины нестабильности результатов и позволяют отделить особенности аппаратной платформы от случайных внешних воздействий.

Важнейшей проблемой при работе с любыми перечисленными метриками является вариативность измерений. Она обусловлена прежде всего тем, что современные операционные системы работают в многозадачном режиме и постоянно перераспределяют ресурсы между конкурентными задачами. Даже при неизменной конфигурации оборудования последовательные измерения могут различаться из-за миграции потоков между ядрами, обслуживания прерываний, подкачки памяти, фоновой индексации файлов, сетевой активности, изменений температурного режима или динамического регулирования частоты процессора. Следовательно, единичное измерение не всегда отражает типичное поведение платформы.

Дополнительным источником вариативности являются эффекты прогрева и предварительного состояния системы. Первые итерации теста могут выполняться в условиях холодных кэшей, неинициализированных структур данных и ещё не стабилизировавшихся частотных режимов. В последующих итерациях результаты зачастую улучшаются или, напротив, ухудшаются, если сказывается нагрев и ограничения энергопотребления. Кроме того, при измерении коротких операций существенную роль может играть накладная стоимость самой инфраструктуры тестирования, поэтому методика должна учитывать достаточное время прогона и осмысленное число повторений.

В связи с этим в данной работе используется подход, при котором измерения организуются как серия испытаний с предварительным разогревом и последующей статистической обработкой. Разогрев позволяет уменьшить влияние неустоявшегося начального состояния. Повторные прогоны дают материал для оценки разброса. Фильтрация выбросов снижает искажение результатов редкими внешними событиями, а использование устойчивых статистических характеристик делает итоговые значения более надёжными. Таким образом, анализ метрик в настоящем исследовании строится не на единичных наблюдениях, а на совокупности данных, интерпретируемых в контексте состояния всей системы.

\section{Обзор существующих инструментов и подходов}

Практика анализа вычислительных систем традиционно основывается на сочетании нескольких классов программных средств, каждое из которых решает собственную задачу. В общем виде такие средства можно разделить на встроенные системные мониторы, синтетические бенчмарки общего назначения, специализированные тесты отдельных подсистем и профилировщики, ориентированные на разработчика. Несмотря на их полезность, ни один из этих классов в отдельности не обеспечивает одновременно полноту контекста, воспроизводимость измерений и удобство прикладной интерпретации результатов.

Первую группу составляют \textbf{встроенные средства операционной системы}. К ним относятся диспетчеры задач, мониторы ресурсов, системные панели наблюдения и консольные утилиты, показывающие загрузку CPU, использование памяти, активность дисков и сети. Преимуществом таких средств является их доступность, оперативность и непосредственная привязка к текущему состоянию платформы. Они незаменимы для диагностики уже наблюдаемой проблемы: например, при резком росте загрузки процессора, нехватке памяти или аномальной дисковой активности. Однако данные инструменты, как правило, ориентированы на моментальный срез состояния и не предназначены для выполнения воспроизводимых серий измерений. Они не формируют контролируемую нагрузку и не дают единого механизма сопоставления разных классов показателей.

Ограниченность встроенных средств особенно заметна в тех случаях, когда требуется не просто зафиксировать факт высокой нагрузки, а ответить на вопрос о предельных возможностях системы. Например, монитор ресурсов может показать текущее потребление CPU или RAM, но не определит пропускную способность памяти при потоковой обработке данных, не оценит латентность случайного доступа и не покажет, каковы возможности платформы в криптографических вычислениях. Кроме того, встроенные инструменты редко дают пользователю научно или методически оформленную интерпретацию результатов, оставляя выводы на усмотрение оператора.

Вторую группу образуют \textbf{синтетические бенчмарки общего назначения}. Наиболее характерный пример такого подхода представляют популярные кросс-платформенные тестовые пакеты, ориентированные на быстрое сравнение систем по набору стандартных нагрузок. Их основное достоинство заключается в удобстве: пользователь получает сводный результат или набор оценок, на основании которых можно сопоставлять разные устройства между собой. Подобные инструменты особенно востребованы в прикладной практике, где необходимо быстро оценить относительный уровень производительности без развёрнутой аналитики.

Однако универсальные бенчмарки общего назначения обладают и существенными ограничениями. Во-первых, они часто функционируют как «чёрный ящик»: пользователь видит итоговый балл или обобщённый набор чисел, но не всегда имеет доступ к деталям методики, параметрам теста и механизму интерпретации. Во-вторых, такие средства, как правило, слабо учитывают текущее состояние системы в момент измерения. Они сообщают результат, но не объясняют, в какой степени на него повлияли фоновые процессы, загрузка диска, нестабильность частоты или иные внешние факторы. В-третьих, универсальные пакеты не всегда позволяют адаптировать тестирование под прикладной профиль пользователя, что ограничивает их ценность при решении конкретных инженерных задач.

Третью группу составляют \textbf{специализированные бенчмарки отдельных подсистем}. Сюда относятся средства оценки производительности процессора, памяти, дисков, сети или конкретных криптографических алгоритмов. Их сильной стороной является фокусировка на одном классе задач, что позволяет реализовать более точную и содержательную методику измерения. Например, тесты памяти могут отдельно оценивать пропускную способность и задержки, дисковые тесты --- различать последовательный и случайный доступ, а криптографические утилиты --- измерять скорость конкретного алгоритма в конкретном режиме работы.

Недостаток специализированных средств состоит в их фрагментарности. Пользователь вынужден обращаться к нескольким независимым программам, каждая из которых использует собственный формат вывода, собственный набор параметров и собственную методику представления результатов. Это затрудняет сопоставление показателей между собой и усложняет формирование целостной картины состояния платформы. Кроме того, раздельное использование множества инструментов не всегда позволяет связать результат отдельного теста с общим состоянием системы в момент выполнения измерения.

Четвёртую группу представляют \textbf{профилировщики и инструменты разработчика}, например семейство \texttt{pprof} в экосистеме Go. Эти средства ориентированы прежде всего на анализ поведения конкретного программного продукта: выявление горячих участков кода, оценку распределения времени выполнения, анализ использования памяти, блокировок и активности горутин. Их значение для оптимизации прикладного ПО чрезвычайно велико, поскольку они позволяют установить, какие именно функции и структуры данных формируют основную нагрузку.

Тем не менее профилировщики не решают задачу первичного системного анализа платформы как таковой. Они требуют наличия исследуемого приложения и репрезентативного сценария его выполнения. Получаемые результаты носят прикладной, а не платформенный характер: например, высокий процент времени в определённой функции объясняет поведение данной программы, но не даёт полноценного ответа на вопрос о предельной пропускной способности памяти системы в целом или о сравнительной эффективности разных аппаратных платформ без привязки к конкретному приложению.

С точки зрения задач компьютерной безопасности следует также учитывать наличие \textbf{криптографически ориентированных средств тестирования}. Они полезны для измерения скорости хеширования, шифрования, генерации ключей и иных операций, однако в большинстве случаев остаются узкоспециализированными утилитами. Такие средства позволяют оценить конкретные алгоритмы, но, как правило, не связывают криптографические показатели с общей телеметрией системы и не встраиваются в более широкий контекст анализа вычислительной платформы.

Обобщая рассмотренные классы инструментов, можно сделать вывод, что существующие подходы либо хорошо решают задачу наблюдения, либо хорошо решают задачу измерения, но редко объединяют их в единой методически целостной форме. Встроенные мониторы дают контекст, но не обеспечивают контролируемого тестирования. Универсальные бенчмарки дают быстрый результат, но часто недостаточно прозрачны и слабо связаны с текущим состоянием системы. Специализированные тесты точны, но фрагментарны. Профилировщики глубоки, но ориентированы на отдельные программы, а не на платформу как целое.

Именно поэтому в рамках настоящей работы обосновывается подход «единое окно», в котором мониторинг, контекстная телеметрия и синтетические бенчмарки интегрируются в одном программном комплексе. С научно-методической точки зрения это позволяет перейти от набора разрозненных показателей к согласованной модели анализа. С практической точки зрения такой подход упрощает сравнение платформ, повышает интерпретируемость результатов и делает инструмент полезным как для инженерной диагностики, так и для учебного применения.

Следовательно, разработка собственного программного комплекса системного анализа представляется обоснованной не как дублирование уже существующих средств, а как попытка объединить их сильные стороны и устранить наиболее существенные ограничения. В дальнейшем это позволяет сформулировать требования к архитектуре системы, способу представления данных и методике выполнения измерений, что и рассматривается в последующих главах работы.
